\chapter{Bảo mật: JWT, OAuth2 và biên tin cậy}
\label{ch:security}

Bảo mật là nơi dự án này dạy được nhiều nhất, vì nó dùng ba cơ chế mà
ngành dùng liên tục --- JWT phi trạng thái, đăng nhập OAuth2, xác thực
hai lớp TOTP --- và nối chúng qua một gateway theo mẫu ``biên tin cậy''
(trusted edge) chuẩn mực.

\section{Spring Security trong một trang}

Spring Security là một chuỗi filter servlet đứng trước các controller
của bạn. Mỗi request đi qua chuỗi đó; các filter xác thực (bạn là ai?)
và các quy tắc phân quyền quyết định (bạn được phép không?). Chuỗi của
auth-service trong dự án, từ \code{SecurityConfig} của nó:

\begin{lstlisting}[language=Java]
http
  .csrf(AbstractHttpConfigurer::disable)                    (*@\co{1}@*)
  .sessionManagement(s ->
      s.sessionCreationPolicy(SessionCreationPolicy.STATELESS)) (*@\co{2}@*)
  .authorizeHttpRequests(auth -> auth
      .requestMatchers("/register", "/login", "/activate",
          "/refresh", "/verify-2fa", "/change-password",
          "/oauth2/**", "/login/oauth2/**",
          "/v3/api-docs/**", "/swagger-ui/**").permitAll()  (*@\co{3}@*)
      .anyRequest().authenticated())
  .oauth2Login(o -> o.successHandler(oAuth2LoginSuccessHandler))
  .addFilterBefore(gatewayHeaderAuthFilter,
      UsernamePasswordAuthenticationFilter.class);          (*@\co{4}@*)
\end{lstlisting}

\begin{description}
  \item[\co{1}] Bảo vệ CSRF che chắn cho phiên cookie của trình duyệt;
  một API JWT phi trạng thái không có phiên cookie nào để lợi dụng, nên
  nó bị tắt. Tắt vì lý do đó --- không phải vì một tutorial bảo thế.
  \item[\co{2}] \code{STATELESS}: không tạo HTTP session. Mỗi request
  phải tự mang bằng chứng danh tính. Đây là tính chất cho phép bất kỳ
  replica nào trả lời bất kỳ request nào --- điều kiện tiên quyết để mở
  rộng.
  \item[\co{3}] Các đường dẫn công khai. Để ý đó là các đường dẫn
  \emph{đã bỏ tiền tố} (\code{/login}, không phải \code{/api/auth/login})
  --- Chương~4 giải thích vì sao. Cũng để ý thứ \emph{vắng mặt}:
  \code{/actuator/**}. Sự vắng mặt đó quyết định chiến lược probe
  Kubernetes của service này ở Chương~17.
  \item[\co{4}] Filter tùy biến tin tưởng các header danh tính từ
  gateway --- chủ đề của Mục~\ref{sec:trusted-edge}.
\end{description}

\section{JWT: danh tính phi trạng thái}

Một JSON Web Token gồm ba phần base64url,
\code{header.payload.signature}. Payload mang các \emph{claim}: subject
(id người dùng), vai trò, hạn dùng. Chữ ký là HMAC-SHA256 trên header
và payload với một secret dùng chung --- chuỗi tối thiểu 256 bit mà cả
auth-service lẫn gateway đều giữ.

Các tính chất quan trọng:

\begin{itemize}
  \item \textbf{Tự chứa.} Xác minh không cần gọi cơ sở dữ liệu: kiểm tra
  chữ ký, kiểm tra hạn, đọc claim. Chính điều này cho phép
  \emph{gateway} xác thực mà không sở hữu dữ liệu người dùng.
  \item \textbf{Đọc được.} Base64 là encoding, không phải mã hóa
  (encryption) --- ai cũng giải mã được payload. Đừng bao giờ đặt bí mật
  trong claim.
  \item \textbf{Không thu hồi được cho tới khi hết hạn.} Token đã ký thì
  không thể ``hủy ký''. Vì thế mới có mẫu hai token: một \emph{access
  token} sống ngắn (ở đây 15 phút) đi cùng một \emph{refresh token} sống
  dài (7 ngày) mà \emph{có} được đối chiếu với cơ sở dữ liệu tại
  \code{/refresh} --- việc thu hồi diễn ra ở đó.
\end{itemize}

\section{Mẫu biên tin cậy}
\label{sec:trusted-edge}

\code{JwtAuthenticationFilter} của gateway xác minh token, rồi chuyển
tiếp request với danh tính dưới dạng header thuần:

\begin{lstlisting}[language=Java]
ServerHttpRequest mutatedRequest = exchange.getRequest().mutate()
    .header(JwtConstants.HEADER_USER_ID, userId)
    .header(JwtConstants.HEADER_USER_ROLE, role)
    .build();
\end{lstlisting}

Ở phía sau, \code{GatewayHeaderAuthFilter} của mỗi service đọc các
header đó và dựng Spring Security context từ chúng --- không xác minh
lại. Service giữ được sự đơn giản; mật mã chỉ chạy một lần.

Nghĩa vụ đi kèm mẫu này: các header đó chỉ đáng tin nếu \emph{mọi}
đường tới một service đều đi qua gateway. Một client với được thẳng tới
course-service có thể giả mạo \code{X-User-Id}. Đây là lý do sâu xa
khiến Ingress \code{/courses} trực tiếp bị xóa (Chương~4), và là lý do
lưu lượng nội cụm vẫn là điểm yếu còn lại --- pod nào cũng vẫn có thể
gọi \code{course-service:8082} với header bịa ra. NetworkPolicy hay
mTLS (service mesh) là cách gia cố trong công nghiệp; Chương~27 bàn khi
nào chúng xứng đáng với độ phức tạp của mình.

\section{Đăng nhập OAuth2 với Google}

Luồng \emph{authorization code} của OAuth2, như nó chạy ở đây:

\begin{enumerate}
  \item Trình duyệt gọi \code{/oauth2/authorization/google} (qua
  gateway); Spring chuyển hướng tới Google kèm client-id của ứng dụng và
  một \code{redirect\_uri}.
  \item Người dùng đồng ý tại Google; Google chuyển hướng trình duyệt về
  \code{redirect\_uri} --- \code{/login/oauth2/code/google} --- mang theo
  một \emph{code} dùng một lần.
  \item Auth-service đổi code + client-secret lấy token theo đường
  server-tới-server, đọc e-mail/hồ sơ người dùng, rồi
  \code{OAuth2LoginSuccessHandler} của dự án tiếp quản: tìm hoặc tạo tài
  khoản cục bộ và phát hành JWT \emph{của chính chúng ta}. Google xác
  thực; ứng dụng vẫn phân quyền bằng token của riêng mình.
\end{enumerate}

\begin{pitfall}[redirect\_uri\_mismatch]
Triệu chứng: Google hiện ``Error 400: redirect\_uri\_mismatch'' thay vì
màn hình đồng ý. Nguyên nhân: URI ứng dụng gửi đi không trùng từng byte
với URI đã đăng ký trong Google Cloud console. Nó thay đổi theo môi
trường (\code{localhost:8080} khi dev, \code{ts.local} trong cụm --- đó
là lý do Deployment ghi đè \code{...GOOGLE\_REDIRECT\_URI}), và mọi
biến thể đều phải được đăng ký. URI phải tới auth-service \emph{chưa bị
bỏ tiền tố}, đó là lý do các route OAuth2 trên gateway không có
StripPrefix.
\end{pitfall}

\section{Xác thực hai lớp TOTP}

Mật khẩu dùng một lần theo thời gian (Time-based One-Time Password,
RFC~6238): khi đăng ký, server sinh một secret, hiện nó dưới dạng mã QR,
và ứng dụng authenticator suy ra mã 6 chữ số từ
\code{HMAC(secret, time/30s)}. Đăng nhập thành hai bước --- kiểm tra mật
khẩu phát ra một trạng thái chờ sống ngắn, rồi \code{/verify-2fa} kiểm
tra mã và phát hành token thật. Server và điện thoại không bao giờ liên
lạc với nhau; chúng chỉ chia sẻ secret và một chiếc đồng hồ.

\section{Ngoài dự án}

Lựa chọn thay thế chính trong công nghiệp đưa việc phát hành token ra
khỏi code của bạn hoàn toàn: một identity provider (Keycloak tự host,
hoặc Auth0/Cognito được quản lý) phát hành token ký bằng cặp khóa
\emph{bất đối xứng} (RS256). Service xác minh bằng khóa \emph{công khai}
lấy từ một URL JWKS --- không có secret dùng chung phải phân phối, xoay
khóa có sẵn:

\begin{lstlisting}[language=yaml]
spring:
  security:
    oauth2:
      resourceserver:
        jwt:
          jwk-set-uri: https://idp.company.com/realms/app/protocol/openid-connect/certs
\end{lstlisting}

Hiểu phiên bản HS256 tự tay làm --- thứ dự án này xây --- chính là điều
làm cho cấu hình kia có ý nghĩa: cùng claim, cùng cách xác minh, chỉ
khác nghi thức khóa.

\begin{quiz}
\begin{enumerate}
  \item Vì sao có thể tắt bảo vệ CSRF cho một API JWT phi trạng thái
  nhưng không thể cho một webapp dùng phiên cookie?
  \item Kẻ tấn công đánh cắp được một access token. Nó dùng được bao
  lâu, trong lúc đó server không thể làm gì, và endpoint nào cuối cùng
  chặn được phiên?
  \item Trong mẫu biên tin cậy, phát biểu chính xác bất biến phải giữ
  để \code{X-User-Id} đáng tin, và nêu hai sai lầm triển khai sẽ phá vỡ
  nó.
  \item Trong luồng Google, secret nào không bao giờ rời server, giá trị
  nào đi qua trình duyệt, và vì sao giá trị trình duyệt thấy được lại an
  toàn?
  \item HS256 so với RS256: ai xác minh được, ai ký được, và vì sao các
  identity provider chọn RS256?
\end{enumerate}
\end{quiz}
